\cotefile{45}
\scanpage{33}\mspage{22}
% archive sheet number: 32
n\textsuperscript{o} 246

la réunion de $]\Delta_1,\Delta_2[$, de $\Delta_1$ et de $\Delta_2$.

\uline{Lorsque} $k$ \uline{satisfait à l'hypothèse du chapeau du n}\textsuperscript{o} 6, on peut, au moyen d'une forme quadratique auxiliaire, établir une correspondance biunivoque \uline{croissante} entre le corps $k$ et l'ensemble ordonné défini par la prop. 6. [N.B. Pour le cas où $k$ ne satisfait pas à l'hypothèse en question, l'adjudant de service fabriquera un contre-exemple et le mettra en exercice]. Pour cela, on introduit dans le plan $E$ n'importe quelle forme quadratique $\phi$ pour laquelle $-\delta$ est un carré, ce qui nous permet d'appliquer tous les résultats des n\textsuperscript{os} précédents. Pour toute demi-droite $\Delta$ autre que $\Delta_{\unsure{o|0}}$, on posera

\[
t(\Delta)=\operatorname{tg}\frac{1}{2}\,(-\Delta_o,\Delta).
\]

Comme on a $\Delta\ne\Delta_o$, $(-\Delta_o,\Delta)$ n'est pas l'angle plat, donc $\frac{1}{2}\,(-\Delta_o,\Delta)$ n'est pas un angle droit et par suite $t(\Delta)\ne\infty$. Il est immédiat que $\Delta\longrightarrow t(\Delta)$ est une bijection de l'ensemble des demi-droites $\ne\Delta_o$ sur $k$. On va montrer que cette bijection est croissante. Soient en effet $\Delta_1,\Delta_2$ deux demi-droites distinctes de $\Delta_o$ ; soit $t_i=t(\Delta_i)\quad(i=1,2)$\unsure{.} Soit $a$ un point, autre que $0$, sur $-\Delta_o$ ; pour la base $(a,Ja)$ de $E$, les demi-droites $\Delta_o,\Delta_1,\Delta_2$ contiendront respectivement (en vertu de la formule (12)) les points :

\[
\begin{aligned}
x_o&=(-1,0)\,,\\
x_i&=(1-t_i^2,\,2t_{\unsure{i|1}})\quad(i=1,2),
\end{aligned}
\]

Pour que la suite $\Delta_o,\Delta_1,\Delta_2$ soit directe, il faut et il suffit que deux au moins des quantités

\[
-2t_1,\quad(1-t_1^2)2t_2-(1-t_2^2)2t_1\,,\quad 2t_2
\]

soient $>0$. Or la seconde s'écrit $2(t_2-t_1)\,(1+t_1t_2)$.
